%% =====================================================================
%%  T-14-08  The Calibrated Question
%%  -------------------------------------------------------------------
%%  One idea per file. Core is mandatory. Slots in canonical order.
%%  Max 700 words. No layout commands. No filler. New source material
%%  for this entry goes in sources/<BOOK>/T-14-08.tex -- never by
%%  rewriting this file (docs/RULES.md R0.1).
%% =====================================================================
%% @kind: topic
%% @id: T-14-08
%% @title: The Calibrated Question
%% @subtitle: Make the other side solve your problem
%% @chapter: c14
%% @order: 80
%% @status: stable
%% @sources: B5
%% @updated: 2026-09-19

\topic{T-14-08}{The Calibrated Question}{Make the other side solve your problem}

\begin{core}
Ask open \emph{how} and \emph{what} questions that hand the counterpart your problem as if it were theirs --- \emph{how am I supposed to do that?}, \emph{what about this works for you?} --- and never \emph{why}, which accuses. The question removes aggression while moving the frame: the counterpart answers, and in answering starts solving on your terms, feeling in control the whole time.
\end{core}
\begin{mechanism}
A calibrated question is an offer of agency with the direction built in. Open-ended how/what questions engage the counterpart in generating the solution, so the outcome is partly authored by them --- the illusion of control is the point: people defend what they feel they chose. Demands trigger resistance; the same content as a question triggers collaboration, which is why the source calls it a way to introduce ideas and requests without sounding pushy --- a nudge that carries no visible push. \emph{Why} is excluded because it assigns blame regardless of intent. The polite-impossible form --- \emph{how am I supposed to do that?} --- is a refusal in the grammar of cooperation: it forces the counterpart to solve your impossibility or lower their demand.
\end{mechanism}
\begin{conditions}
Strongest against aggressive or positional counterparts, in asymmetric-power talks where a direct no is expensive, and mid-conflict where a demand would escalate. Weakens with counterparts who will not think --- the question needs an interlocutor --- and under time pressure where a decision, not a dialogue, is required; overuse turns every exchange into therapy and the counterpart notices they are being steered by questions.
\end{conditions}
\begin{application}
  \item Convert your next demand into \emph{how} or \emph{what}: not \emph{you need to lower the price} but \emph{how am I supposed to take this to my side?}
  \item Map the field before answering theirs: \emph{what is your objective here?} \emph{what about this is important to you?}
  \item For the impossible ask, use the polite-impossible: \emph{how am I supposed to do that?} --- sincere tone, genuine puzzle.
  \item Stack questions; never stack demands. Each answer should pull the next question into reach.
  \item Save \emph{why} for genuine curiosity about the past, never for their present behaviour.
\end{application}
\begin{feedback}
  \item Landing: the counterpart starts designing within your constraint --- their answer contains your terms.
  \item Landing: tone de-escalates measurably; they begin asking \emph{we} questions.
  \item Failing: \emph{just figure it out} --- they decline the authorship; fall back to labels (T-14-07) before re-asking.
  \item Failing: they answer your how-question with a question of their own; two steering systems are now running --- decide whose field you are on.
\end{feedback}
\begin{failure}
Detected, the pattern reads as manipulation-by-questions, and the counterpart feels managed exactly when they believed they were choosing --- the backlash is colder than a refused demand. It also fails as deflection: a question that only stalls, without accepting any answer, is visible within two rounds. And chronic use atrophies the direct ask; some moments need a stated position, and a person who only ever questions becomes a person who cannot be negotiated with.
\end{failure}
\begin{countermeasures}
  \item Notice you are being asked how \emph{you} will solve \emph{their} problem; answer once, then ask what they are bringing.
  \item Return the question with its frame exposed: \emph{that's a fair question --- here's my constraint; what's yours?}
  \item Refuse the authorship trap: do not design inside their constraint. \emph{That's your number to solve} is a complete answer.
  \item Track why-questions aimed at your present behaviour; they are accusations wearing grammar (T-06-04).
\end{countermeasures}

\sources{B5}
\seesources
\seealso{T-14-01, T-16-01, T-24-02}
